% Supplementary Information of the ERL Letter (D-107): the red team, the three strongest reasons
% the headline could be wrong, each tested (G3 memo section 9; D-082, D-096; red-team runs).
% Every number is a \metric{} macro.

Before the venue decision, the headline was attacked on three counts: that the ``observed'' shorelines, not the grids, are wrong (R1); that a datum varying across the shelf, not the shape of the seabed, causes the errors (R2); and that the near-term correction cannot be trusted below the observed range (R3).

\subsection*{R1: the observed shorelines are the error}
If the observed extents were biased by turbid or dark water, wet salt flats, reeds, wind setup, the record-low rule or the rules for cells outside the observed cover, the grids would look wrong when the observations are. Five tests address this. (i) The noise floor from two independent halves of the observations is built into the decision rule, and the grids exceed it at most levels. (ii) The hindcast was rerun without the record-low rule, with the shape constants refitted (table~\ref{tab:si_r1r2}). (iii) The audit on E2 and E3 only, without the rising limb, keeps every grid's majority outcome (\metric{erl.text.robust.e2e3.majority.n_sectors_min} sectors for each of \metric{erl.text.h1.majority_products}). (iv) The threshold, sensor, season and wind-setup variants change sector-level details and the elevation model's magnitudes (section~\SIrobust), but they were judged on the elevation model, not on the audit. (v) At Great Salt Lake the same observation pipeline passes a survey-quality surface (section~\SIgsl).

Without the record-low rule only the open shelf changes: there, deep and clear water is misread as land in a share of acquisitions, which inflates the open-shelf noise floor and displacements alike. GEBCO\_2026's open-shelf median is \metric{redteam.noenv.shape.gebco2026.open_shelf_km}~km without the rule against \metric{redteam.primary.shape.gebco2026.open_shelf_km}~km with it. Because the open shelf is in the comparison group, the deltaic and vegetated contrast loses significance for every grid. R1 does not sink the first half of the audit hypothesis; it weakens the second.

\subsection*{R2: a spatially varying datum causes the errors}
The shape variant removes one constant per grid. A level reference that varies across the shelf (setdown in Komsomolets Bay, geoid-model differences, or a mismatch between DAHITI and the grids that varies in space) would look like shape error. The constant was therefore refitted per grid and sector on training cells, smoothed, and the hindcast rerun. The sector constants span \metric{redteam.sector_shape.offset_min.gebco2026_m} to \metric{redteam.sector_shape.offset_max.gebco2026_m}~m for GEBCO\_2026 and \metric{redteam.sector_shape.offset_min.etopo2022_m} to \metric{redteam.sector_shape.offset_max.etopo2022_m}~m for ETOPO 2022. Even so, the grids still mispredict in most sectors (table~\ref{tab:si_r1r2}); the only sector that changes for the GEBCO grids is Komsomolets Bay, whose level departs from the basin's, and the deltaic and vegetated contrast becomes stronger. R2 does not explain the errors.

\begin{table}[!htbp]
\centering
\small
\caption{\textbf{Red-team tests R1 and R2 of the hindcast audit.} Sectors (of 7) in which each grid mispredicts after its vertical constant, and the one-sided $p$ of the block test of larger errors in the deltaic and vegetated sectors: in the reported run, without the record-low rule (R1), and with a constant per sector (R2).}
\label{tab:si_r1r2}
\begin{tabular}{lrrrrrr}
\hline
 & \multicolumn{2}{c}{Reported run} & \multicolumn{2}{c}{Without record-low rule} & \multicolumn{2}{c}{Constant per sector} \\
Grid & Sectors & $p$ & Sectors & $p$ & Sectors & $p$ \\
\hline
GEBCO\_2026 & \metric{h1.mispredicts.shape.gebco2026.n_sectors} & \metric{erl.text.p4.h1.block_test.shape.gebco2026.p} & \metric{redteam.noenv.shape.gebco2026.n_sectors} & \metric{erl.text.p4.redteam.noenv.shape.gebco2026.block_p} & \metric{redteam.sector_shape.sector_shape.gebco2026.n_sectors} & \metric{erl.text.p4.redteam.sector_shape.sector_shape.gebco2026.block_p} \\
GEBCO\_2022 & \metric{h1.mispredicts.shape.gebco2022.n_sectors} & \metric{erl.text.p4.h1.block_test.shape.gebco2022.p} & \metric{redteam.noenv.shape.gebco2022.n_sectors} & \metric{erl.text.p4.redteam.noenv.shape.gebco2022.block_p} & \metric{redteam.sector_shape.sector_shape.gebco2022.n_sectors} & \metric{erl.text.p4.redteam.sector_shape.sector_shape.gebco2022.block_p} \\
ETOPO 2022 & \metric{h1.mispredicts.shape.etopo2022.n_sectors} & \metric{erl.text.p4.h1.block_test.shape.etopo2022.p} & \metric{redteam.noenv.shape.etopo2022.n_sectors} & \metric{erl.text.p4.redteam.noenv.shape.etopo2022.block_p} & \metric{redteam.sector_shape.sector_shape.etopo2022.n_sectors} & \metric{erl.text.p4.redteam.sector_shape.sector_shape.etopo2022.block_p} \\
ETOPO1 & \metric{h1.mispredicts.shape.etopo1.n_sectors} & \metric{erl.text.p4.h1.block_test.shape.etopo1.p} & \metric{redteam.noenv.shape.etopo1.n_sectors} & \metric{erl.text.p4.redteam.noenv.shape.etopo1.block_p} & \metric{redteam.sector_shape.sector_shape.etopo1.n_sectors} & \metric{erl.text.p4.redteam.sector_shape.sector_shape.etopo1.block_p} \\
SRTM15+ & \metric{h1.mispredicts.shape.srtm15p.n_sectors} & \metric{erl.text.p4.h1.block_test.shape.srtm15p.p} & \metric{redteam.noenv.shape.srtm15p.n_sectors} & \metric{erl.text.p4.redteam.noenv.shape.srtm15p.block_p} & \metric{redteam.sector_shape.sector_shape.srtm15p.n_sectors} & \metric{erl.text.p4.redteam.sector_shape.sector_shape.srtm15p.block_p} \\
\hline
\end{tabular}
\end{table}

\subsection*{R3: the near-term correction cannot be trusted below the observed range}
The registered result, that correction lowers emergence at \metric{h3.scenario.1_m}~m by a factor of \metric{erl.text.h3.registered.1.unc_over_fused_min}--\metric{erl.text.h3.registered.1.unc_over_fused_max}, rests on an error model fitted at observed levels and applied below them. Four findings narrow it (section~\SInearterm). (i) The leave-one-epoch-out test fails at the lowest epoch: fitted without E3 and tested on it, the corrected surface's emergence error exceeds the uncorrected grid's for \metric{erl.text.h3.registered.e3_n_corrected_worse} of \metric{erl.text.h1.n_products} grids. (ii) The error model is degenerate in many sectors: in \metric{h3.atl24.error_model.n_outside_fit_range} of \metric{h3.atl24.error_model.n_models} grid--sector models the shift applied in the near-term zone lies outside the shifts the model was fitted on, reaching \metric{h3.atl24.error_model.max_applied_shift_m}~m against fitted shifts of \metric{h3.atl24.error_model.fit_shift_min_m} to \metric{h3.atl24.error_model.fit_shift_max_m}~m. Whole-metre depths put many units at nearly the same model depth, so the Theil--Sen slopes between them are steep; the hold stops extrapolation in depth but not in shift, and the registered guard does not catch steep positive slopes. Near-term emergence then collapses in those sectors, which accounts for much of the factor. (iii) No independent seabed check was possible: ATL24 holds no data over the Caspian. (iv) Under the registered constant-form sensitivity, which stays inside the fitted shifts by construction, the change at \metric{h3.scenario.1_m}~m is no longer decision-relevant for \metric{erl.text.h3.constant.not_relevant_1m.n_products} of \metric{erl.text.h1.n_products} grids, and it too fails at the lowest epoch (\metric{erl.text.h3.constant.e3_n_corrected_worse} of \metric{erl.text.h1.n_products} grids). R3 partly sinks the near-term hypothesis as worded: the registered verdict stands as registered, but its magnitude is an artefact and no correction is validated where it is applied. The Letter therefore states only what survives without any error model.
